% ATTEMPT_STATUS: unresolved
% TOP_PROBLEM: {"problemId": "problem.finiteness-of-central-configurations-for-the-n-body-problem-n-4", "canonicalTitle": "Finiteness of central configurations for the n-body problem, n >= 4", "releaseRank": 410, "primaryDomain": "mathematical_physics", "primaryDomainLabel": "Mathematical physics", "displayStatus": "open_with_solved_subcases", "releaseStatus": "open_with_solved_subcases"}
% !TeX program = lualatex
\documentclass[11pt]{article}
\usepackage[margin=25mm]{geometry}
\usepackage{fontspec}
\setmainfont{Latin Modern Roman}
\usepackage{amsmath,amssymb,mathtools}
\usepackage{unicode-math}
\setmathfont{Latin Modern Math}
\usepackage{xurl,hyperref}
\hypersetup{colorlinks=true,urlcolor=blue}
\setcounter{secnumdepth}{0}
\setlength{\parindent}{0pt}
\setlength{\parskip}{5pt}
\setlength{\emergencystretch}{3em}
\begin{document}
\section*{\texorpdfstring{Finiteness of central configurations for the n-body problem, n >= 4}{Finiteness of central configurations for the n-body problem, n >= 4}}\label{410}
\subsection{Definitions and mathematical statement}
For positive masses $m_1,\ldots,m_n$, a collision-free Newtonian configuration $q_1,\ldots,q_n\in{\mathbb{R}}^d$ is central if, with center of mass $c=\sum_i m_iq_i/\sum_i m_i$, there is $\lambda>0$ such that
\[
 \sum_{j\ne i}m_j\frac{q_j-q_i}{\|q_j-q_i\|^3}
 =-\lambda(q_i-c)\qquad(1\le i\le n).
\]
The standard finiteness question counts configurations modulo translation, rotation, and positive rescaling, equivalently after $c=0$ and $\sum_i m_i|q_i|^2=1$, modulo orthogonal transformations. It asks for finiteness for every fixed mass vector and $n\ge4$. Without this quotient the literal count is infinite by scaling. The catalogue and its short AIM question do not fix the planar versus spatial version; $d$ must be specified in applying the question, rather than silently conflating these configuration spaces.
\par\smallskip\textit{Formulation and concept references:} \href{https://aimath.org/WWN/varcelest/varcelest.pdf}{[S1]}.

\subsection{Short English statement}
For fixed positive masses, are there only finitely many shapes of Newtonian central configurations after removing translation, rotation, and scaling freedoms?
\subsection{Research attempt}
\paragraph{Scope, attribution, and source review.}
Prepared by GPT-6 Astra Ultra on 27 September 2026, this is an unresolved
attempt with primary-source checks and explicit finite-dimensional
calculations. Fix the positive mass vector and specify the ambient
dimension $d$. Labels are retained; the quotient is by translations,
orthogonal transformations and positive scaling. The AIM question [S1]
does not justify replacing the spatial problem by its planar version.
Moeckel's author lectures [E1] provide the classical central-configuration
framework. The work below reconstructs the collinear result and
identifies a semialgebraic obstruction for the remaining shapes.
Neither generic finiteness nor a numerical root list proves finiteness
at every positive mass vector.

\paragraph{Normalization and the variational equation.}
Set
\[
 U(q)=\sum_{i<j}\frac{m_im_j}{|q_i-q_j|},\qquad
 I(q)=\sum_i m_i|q_i|^2,\qquad \sum_i m_iq_i=0.
\]
The Euclidean gradient is
$\nabla_{q_i}U=m_i\sum_{j\ne i}m_j(q_j-q_i)/|q_j-q_i|^3$.
Hence centrality is $\nabla U=-\lambda(m_iq_i)_i$.
Euler's identity for a homogeneous function of degree $-1$ yields
\[
 -U=\sum_iq_i\cdot\nabla_{q_i}U=-\lambda I,\qquad
                         \lambda=U/I>0.
\]
Thus no additional positive multiplier is arbitrary after $I=1$.
Under $q\mapsto aq$ the multiplier changes to $\lambda/a^3$.
Every central shape has a unique scaling with $\lambda=1$ and
a unique scaling with $I=1$, although these are generally different.

\paragraph{A complete collinear calculation.}
Place the bodies on a line with a fixed order
$x_1<\cdots<x_n$. On that open convex chamber consider
\[
 F(x)=\sum_{i<j}\frac{m_im_j}{x_j-x_i}
                  +\frac12\sum_i m_ix_i^2.
\]
For any real increment $h$,
\begin{equation}\label{a410:hessian}
 D^2F(x)[h,h]
 =2\sum_{i<j}\frac{m_im_j(h_j-h_i)^2}{(x_j-x_i)^3}
                      +\sum_i m_ih_i^2>0
 \quad(h\ne0).
\end{equation}
The positivity of every mass is used here. The function tends to
infinity at a collision and at infinity. More explicitly, a sublevel
bounds all $|x_i|$ by its quadratic term, and bounds each inverse
distance by the corresponding positive potential summand.
Thus its sublevels are compact inside the chamber, and it has a
minimum. Strict convexity makes this minimum the unique critical point.

At a critical point the equations are exactly centrality with
$\lambda=1$. Summing the gradient equations gives
$\sum_i m_ix_i=0$, because internal forces cancel, so centering was
not an extra unsatisfied constraint. Conversely every central
configuration of this order rescales to such a critical point.
There is therefore precisely one collinear shape for each ordering
up to orientation reversal. For labelled bodies the number is
$n!/2$ when $n\geq2$. Equal masses do not change this count unless
one also quotients by label permutations, which is not the convention
used in this calculation. This is the classical Moulton result,
derived here rather than extrapolated to higher dimensions.

\paragraph{Why this convexity proof stops in the plane.}
For one pair with displacement $r=q_i-q_j$ and increment
$h=h_i-h_j$, the second derivative of its potential is
\[
 m_im_j\left(\frac{3(r\cdot h)^2}{|r|^5}
                         -\frac{|h|^2}{|r|^3}\right).
\]
It is positive for longitudinal increments and negative for
nonzero transverse increments. On a line only the first situation
occurs and yields \eqref{a410:hessian}; in dimension at least two,
both occur. The positive quadratic term need not overcome the
negative transverse contribution near a close pair. Therefore
$U+I/2$ is not globally strictly convex on a planar configuration
chamber. Choosing an ordering of $x$-coordinates does not remove
the transverse directions.

Compactness of a normalized critical set, when available, also does
not imply finiteness. A circle is a compact infinite critical set of
a smooth function, for instance of $(x^2+y^2-1)^2$. The missing
ingredient is isolation after the symmetry quotient, not just
control of escape or collisions.

\paragraph{Another complete restricted case: full affine dimension.}
Suppose the $n$ bodies are affinely independent, so they span
dimension $n-1$. Put $M=\sum_i m_i$. The center-of-mass identity is
\[
 q_i-c=-M^{-1}\sum_{j\ne i}m_j(q_j-q_i).
\]
Substitution into centrality gives, for each $i$,
\[
 \sum_{j\ne i}m_j\left(|q_j-q_i|^{-3}-\lambda/M\right)
                                      (q_j-q_i)=0.
\]
The $n-1$ displayed displacement vectors are linearly independent.
Every coefficient must vanish, so every pairwise distance equals
$(M/\lambda)^{1/3}$. Thus the shape is a regular simplex.
Conversely the same identity verifies centrality of a regular
simplex for every positive mass vector, with its mass center
possibly different from its geometric center. Up to the stated
quotient there is exactly one such full-dimensional shape.
This covers the nonplanar four-body simplex shape, but not the
planar four-body configurations or intermediate affine dimensions.

\paragraph{Encoding the actual shape quotient.}
For a fixed $d$, centered configurations can be represented by their
Gram matrix $G_{ij}=q_i\cdot q_j$. Two labelled configurations have
the same Gram matrix if and only if an orthogonal transformation
identifies them in $\mathbb R^d$. The forward implication is
immediate. For the converse, the correspondence between their
spanning vectors preserves all inner products, is well-defined
on their spans, and extends to an orthogonal transformation.

Introduce $t_{ij}=|q_i-q_j|^{-1}>0$ for $i<j$.
The normalized collision-free shape set is described by
\[
 G\succeq0,\quad \operatorname{rank}G\leq d,\quad
 G(m_1,\ldots,m_n)^{\mathsf T}=0,\quad
 \sum_i m_iG_{ii}=1,\quad
 t_{ij}^2(G_{ii}+G_{jj}-2G_{ij})=1.
\]
Positive semidefiniteness and the rank bound are polynomial
inequalities and equations on minors. Centrality can also be
encoded without selecting coordinates: require, for every $i,k$,
\begin{equation}\label{a410:gram}
 \sum_{j\ne i}m_jt_{ij}^3(G_{jk}-G_{ik})+\lambda G_{ik}=0,
 \qquad \lambda=\sum_{i<j}m_im_jt_{ij}.
\end{equation}
Indeed the left side is the inner product of the centrality
residual with $q_k$. That residual lies in the span of all $q_k$;
orthogonality to every spanning vector forces it to vanish.
The final multiplier equation is $I=1$ and Euler's identity.
Conversely a positive semidefinite Gram matrix of the required
rank is realized by vectors, so the encoding loses no shapes.

The resulting set is semialgebraic, with the fixed real masses as
coefficients. A zero-dimensional semialgebraic set is finite; this
follows from its finite cell decomposition into points. Therefore
an infinite family at one mass vector would contain a
positive-dimensional semialgebraic cell and hence a nonconstant
smooth arc of normalized Gram matrices. This is a concrete
reduction: the sought obstruction is a genuine continuous
family after rotation and scaling have already been removed.

\paragraph{What degeneracy tests can and cannot prove.}
On a smooth coordinate chart for shapes, an invertible Jacobian of
the central-configuration equations isolates a solution by the
implicit function theorem. A compact collection of isolated points
is finite if it has no additional accumulation points in the
critical set. More generally the semialgebraic reduction needs no
compactness once dimension zero is proved. But showing that a
Jacobian is invertible at a computed list neither proves that the
list is exhaustive nor excludes degenerate solutions elsewhere.

A generic-mass argument has a second quantifier issue. The toy
system $ax=0$, $y=0$ has one solution for every $a\ne0$ and an entire
line at $a=0$. A positive mass vector can lie on an algebraic
exceptional locus, such as one defined by a relation among its
positive coordinates. This toy system is not a central
configuration; it demonstrates the logical gap in passing from
generic finiteness to every mass vector.

\paragraph{Verification and final status.}
The force sign, scale exponent, Hessian and Gram residual equations
were checked symbolically; the simplex test verifies the multiplier
normalization for unequal masses. All arguments keep collisions
excluded and labels fixed. The remaining general step is to rule
out positive-dimensional semialgebraic families for every fixed
positive mass vector in the specified $d$, or to construct one.
\textbf{UNRESOLVED.} The collinear and full-affine-dimensional cases
are complete, while the general shape space is reduced to a
specific algebraic isolation problem. No new finiteness theorem
for all planar or spatial configurations is claimed.

\subsection{Sources}
\begin{itemize}
\item[S1]American Institute of Mathematics, Variational Methods in Celestial Mechanics problem list, Problem 24 (2004).
\url{https://aimath.org/WWN/varcelest/varcelest.pdf}.
\textit{Formulation source. Consulted 24 September 2026: Central-configuration finiteness problem and its stated scope.}
\item[E1]Richard Moeckel, \emph{Lectures on Central Configurations}, author notes, 20 March 2014.
\url{https://www-users.cse.umn.edu/~rmoeckel/notes/CentralConfigurations.pdf}.
\textit{Primary author account of the variational and dimensional framework; consulted 27 September 2026. The collinear and simplex deductions are given in full above.}
\end{itemize}
\end{document}
